% Esquisse d'un programme (cote 119, copy 1): archive cote 119, site page 19 (PDF page 20), typescript page 17.
% typescript. Draft compared: EsquisseFr.pdf (text layer):625-694.
% AI-transcribed by Claude Opus 5.5 (claude-opus-5-5, adaptive thinking, effort high), Message Batches API, batch pilot (msgbatch_0189Nhf1xB615YWnztubtLfc), 2026-11-12; not yet checked by eye.
\scanpage{119/19}
\mspage{17}

sous une multiplicité d'avatars (dont certains ont été effleurés
dans le présent rapport), et qui du point de vue de la théorie de
Galois-Teichmüller peut être considéré comme la ``pierre de cons-%
truction'' fondamentale de la ``tour de Teichmüller''. L'élément de
structure de $S\ell(2,\mathbb{Z})$ qui me fascine avant tout, est bien sûr
l'action extérieure du groupe de Galois $\mathrm{I}\!\Gamma$ sur le compactifié
profini. Par le théorème de Bielyi, prenant les compactifiés pro-%
finis de sous-groupes d'indice fini de $S\ell(2,\mathbb{Z})$, et l'action ex-%
térieure induite (quitte à passer également à un sous-groupe ou-%
vert de $\mathrm{I}\!\Gamma$), \uline{on trouve essentiellement les groupes fondamentaux}
\uline{de toutes les courbes algébriques} (pas nécessairement compactes)
\uline{définis sur des corps de nombres} $K$, \uline{et l'action extérieure de}
$\mathrm{Gal}(\bar{K}/K)$ \uline{dessus} - du moins est-il vrai que tout tel groupe fon-%
damental apparaît comme quotient d'un des premiers groupes (*).
Tenant compte du ``yoga anabélien'' (qui reste conjectural), disant
qu'une courbe algébrique anabélienne sur un corps de nombres $K$
(extension finie de $\mathbb{Q}$) est connu à isomorphisme près quand on
connaît son groupe fondamental mixte (ou ce qui revient au même,
l'action extérieure de $\mathrm{Gal}(\bar{K}/K)$ sur son groupe fondamental profini
géométrique), on peut donc dire que \uline{toutes les courbes algébriques}
\uline{définies sur des corps de nombres sont} ``\uline{contenues}'' \uline{dans le compac-}%
\uline{tifié profini} $S\ell(2,\mathbb{Z})\hat{}$, \uline{et dans la connaissance d'un certain sous-}%
\uline{groupe} $\mathrm{I}\!\Gamma$ \uline{du groupe des automorphismes extérieurs de ce dernier} !
Passant aux abélianisés des groupes fondamentaux précédents, on voit
notamment que toutes les représentations abéliennes $\ell$-adiques chères
à Tate et consorts, définies par des jacobiennes et jacobiennes gé-%
néralisées de courbes algébriques définies sur des corps de nombres,
sont contenues dans cette seule action de $\mathrm{I}\!\Gamma$ sur le groupe profini
anabélien $S\ell(2,\mathbb{Z})\hat{}$!($^{4}$)

\rule{5cm}{0.4pt}

(*) En fait, il s'agit de quotients de nature particulièrement tri-%
viale, par des sous-groupes abéliens produits de ``modules de Tate''
$\hat{\mathbb{Z}}(1)$, correspondants à des ``groupes-lacets'' autour de points à l'infini.

\margin{19}
